\chapter{Considerações Finais}

Este TCC I apresentou a fundamentação teórica da persistência poliglota, o estado da arte
em importação e modelagem de dados heterogêneos, e a proposta da ferramenta
\textit{Polyglot Import CSV} como utilitário de \textit{bootstrap} declarativo a partir
de CSV.

As principais contribuições desta etapa são: (i) um formato de configuração JSON validado
estaticamente por JSON Schema, separado em mapeamento e conexão, com mapeamento recursivo
de colunas e filtros por valor; (ii) um protótipo em Python que materializa o mesmo CSV
operacional em PostgreSQL, MongoDB, Cassandra, Redis e Neo4j; (iii) o modo
\texttt{-{}-dry-run} para revisão prévia de contagens; e (iv) o cenário e-commerce
reprodutível com \texttt{ecommerce\_join.csv}, \texttt{import\_config.json},
\texttt{sgbd\_config.json} e \texttt{docker-compose.yml}.

As limitações reconhecidas incluem: testes automatizados predominantemente unitários (sem
suíte de integração contínua contra Docker); processamento em memória do CSV inteiro; e
dependência do \textit{driver} Cassandra em ambientes Python recentes. Itens como
interface gráfica, suporte a TSV/Excel, \textit{chunked processing} e novos conectores
permanecem no escopo do TCC II (Capítulo 5).

O trabalho de conclusão de curso em dezembro de 2026 poderá aprofundar a avaliação
experimental (desempenho, consistência entre \textit{stores}) e a interface de uso,
consolidando a ferramenta como artefato de pesquisa e de código aberto alinhado à
literatura sobre SGPD e modelagem poliglota de dados.
